%==================================================================
%   Studienprojekt 2 — Sensor-Based Fall Detection with ML
%   Aly Elbermawy · TH Mittelhessen · 2026
%
%   Compile with:
%     pdflatex main
%     bibtex   main
%     pdflatex main
%     pdflatex main
%==================================================================

\documentclass[11pt,a4paper,oneside]{report}

% -------------- Encoding & language --------------
\usepackage[utf8]{inputenc}
\usepackage[T1]{fontenc}
\usepackage[english]{babel}
% Hinweis: Kapitel 3 ist auf Deutsch verfasst. Fuer korrekte deutsche
% Silbentrennung auf deinem System (TeX Live mit hyphen-german) bitte
% ersetzen durch:
%     \usepackage[ngerman,english]{babel}
% und um jedes deutsche Kapitel \selectlanguage{ngerman} ... \selectlanguage{english} setzen.
% \usepackage{lmodern}  % not available in this environment; default Computer Modern is used

% -------------- Page geometry & layout -----------
\usepackage[a4paper,left=3cm,right=2.5cm,top=2.5cm,bottom=2.5cm]{geometry}
\usepackage{setspace}
\onehalfspacing
\usepackage{parskip}

% -------------- Math -----------------------------
\usepackage{amsmath,amssymb,amsthm}

% -------------- Graphics & floats ----------------
\usepackage{graphicx}
\usepackage{float}
\usepackage{caption}
\captionsetup{font=small,labelfont=bf}

% -------------- Tables ---------------------------
\usepackage{booktabs}
\usepackage{tabularx}
\usepackage{array}
\newcolumntype{L}[1]{>{\raggedright\arraybackslash}p{#1}}
\newcolumntype{C}[1]{>{\centering\arraybackslash}p{#1}}

% -------------- Citations & bibliography ---------
\usepackage[numbers,square,sort&compress]{natbib}
\bibliographystyle{ieeetr}

% -------------- Hyperlinks (loaded last) ---------
\usepackage[hidelinks,pdfusetitle]{hyperref}

% -------------- Nomenclature / acronyms ----------
\usepackage[printonlyused]{acronym}

% -------------- Header/footer --------------------
\usepackage{fancyhdr}
\pagestyle{fancy}
\fancyhf{}
\fancyhead[L]{\small\nouppercase{\leftmark}}
\fancyfoot[C]{\thepage}
\renewcommand{\headrulewidth}{0.4pt}

% -------------- Chapter title style (subtle) -----
\usepackage{titlesec}
\titleformat{\chapter}[hang]
  {\normalfont\huge\bfseries}{\thechapter}{1em}{}
\titlespacing*{\chapter}{0pt}{-20pt}{20pt}

%==================================================================
%   Document metadata
%==================================================================
\title{3-, 6- and 9-Axis Sensor Systems for Fall Detection: A Comparative Study}
\author{Aly Elbermawy}
\date{July 2026}

%==================================================================
\begin{document}

%------------------------------------------------------------------
%   TITLE PAGE
%------------------------------------------------------------------
\begin{titlepage}
\begin{center}

% --- THM logo placeholder (replace with real logo image) ---
% \includegraphics[width=6cm]{thm_logo.pdf}
\vspace*{1cm}
{\LARGE\bfseries TH Mittelhessen\par}
\vspace{0.2cm}
{\small Technische Hochschule Mittelhessen - University of Applied Sciences\par}

\vspace{2.5cm}
{\Large\bfseries Studienprojekt 2\par}
\vspace{0.4cm}
{\large Wissenschaftliches Arbeiten\par}

\vspace{2cm}
\rule{0.7\textwidth}{0.4pt}\\[0.4cm]
{\Large\bfseries Subject:\par}
\vspace{0.3cm}
{\Large\bfseries Sensor-Based Fall Detection with Machine\\Learning}\\[0.4cm]
\rule{0.7\textwidth}{0.4pt}

\vspace{2cm}
\begin{tabbing}
\hspace{4.5cm}\= \kill
\textbf{Author:}         \> Aly Elbermawy \\
\textbf{Matriculation number:} \> 5335848 \\[0.4cm]
\textbf{Supervisor:}     \> Dr.\ Kovalev \\
\textbf{Industry Partner:} \> TrackTech GmbH \\[0.4cm]
\textbf{University:}     \> Technische Hochschule Mittelhessen, Friedberg \\[0.4cm]
\textbf{Submitted on:}   \> July~24, 2026
\end{tabbing}

\end{center}
\end{titlepage}

%------------------------------------------------------------------
%   CONFIDENTIALITY / SPERRVERMERK
%------------------------------------------------------------------
\thispagestyle{empty}
\vspace*{2cm}
\section*{Confidentiality Clause}
Transmittal, reproduction, dissemination and/or editing of this document as well as
utilization of its contents and communication thereof to others without express
authorization are prohibited. Offenders will be held liable for payment of damages.
All rights created by patent grant or registration of a utility model or design
patent are reserved.

\vspace{1.5cm}
\section*{Sperrvermerk}
Die vorliegende Arbeit beinhaltet interne vertrauliche Informationen der TrackTech
GmbH bzw. Clever-Sole. Die Weitergabe des Inhaltes der Arbeit und eventuell
beiliegender Zeichnungen und Daten im Gesamten oder in Teilen ist grundsätzlich
untersagt. Es dürfen keinerlei Kopien oder Abschriften -- auch in digitaler Form
-- gefertigt werden. Ausnahmen bedürfen der schriftlichen Genehmigung der TrackTech
GmbH.
\newpage

%------------------------------------------------------------------
%   ABSTRACT
%------------------------------------------------------------------
\thispagestyle{empty}
\section*{Abstract}
Falls are a leading cause of severe injury and reduced quality of life for elderly
and care-dependent individuals. A reliable, energy-efficient and everyday-wearable
fall-detection system can raise alarms in real time and thereby improve care safety.
The novelty of the present project lies in the target form factor: a
shoe-insole-integrated inertial sensor, rather than the wrist, hip or waist positions
that dominate most public datasets.

The present study, submitted in partial fulfilment of the module \emph{Studienprojekt~2}
at TH Mittelhessen, addresses the first of four assigned tasks. Its objective is a
literature- and datasheet-based comparison of 3-, 6- and 9-axis sensor configurations
in the context of fall detection. The research question is whether a higher number
of sensed axes actually improves detection performance and, if so, up to which point
the improvement justifies the additional cost in energy, computation and complexity.

The analysis draws on peer-reviewed evaluations of accelerometer-only, IMU-based and
MARG-based fall detection algorithms, complemented by manufacturer datasheets of
three representative components (MPU-6050, LSM6DSOX, BHI360). It concludes that the
transition from three to six axes yields a well-documented, substantial improvement
-- primarily in specificity, i.e.\ fewer false alarms -- whereas the further step
to nine axes adds absolute orientation whose benefit is marginal for the
short-duration fall event itself, while introducing new susceptibility to indoor
magnetic disturbances. A 6-axis configuration is therefore recommended as the sweet
spot for the insole use case, which is also compatible with the leading public
6-axis fall datasets (SisFall, KFall) used in subsequent tasks.
\newpage

%------------------------------------------------------------------
%   CONTENTS + LISTS
%------------------------------------------------------------------
\pagenumbering{roman}
\setcounter{page}{3}

\tableofcontents
\newpage

\listoffigures
\listoftables
\newpage

%------------------------------------------------------------------
%   ABBREVIATIONS
%------------------------------------------------------------------
\chapter*{Abbreviations}
\addcontentsline{toc}{chapter}{Abbreviations}

\begin{acronym}[LSM6DSOX]
  \acro{ADL}{Activity of Daily Living}
  \acro{ADC}{Analog-to-Digital Converter}
  \acro{AHRS}{Attitude and Heading Reference System}
  \acro{DoF}{Degrees of Freedom}
  \acro{IMU}{Inertial Measurement Unit}
  \acro{MARG}{Magnetic, Angular Rate and Gravity}
  \acro{MEMS}{Micro-Electro-Mechanical Systems}
  \acro{ML}{Machine Learning}
  \acro{MLC}{Machine Learning Core}
  \acro{ODR}{Output Data Rate}
  \acro{SVM}{Signal Vector Magnitude}
\end{acronym}
\newpage

%==================================================================
\pagenumbering{arabic}
\setcounter{page}{1}

%------------------------------------------------------------------
\chapter{Einleitung}
\label{ch:intro}

\section{Überblick}
Stürze gehören zu den häufigsten Ursachen für schwere Verletzungen und
langfristige gesundheitliche Komplikationen bei älteren und pflegebedürftigen
Menschen, insbesondere bei Personen mit Demenz. Ein zuverlässiges,
energieeffizientes und alltagstaugliches automatisches Sturzerkennungssystem
auf Basis tragbarer Inertialsensoren kann dazu beitragen, Alarme schnell
auszulösen und dadurch die Sicherheit in der Pflege zu erhöhen.

Die Qualität der automatischen Sturzerkennung hängt nicht nur vom verwendeten
Machine-Learning-Modell zur Klassifikation ab, sondern auch von der zugrunde
liegenden Sensorhardware und den verfügbaren Trainingsdaten. Das vorliegende
\emph{Studienprojekt~2} ist daher in vier Aufgaben gegliedert: eine kurze
theoretische Bewertung der Sensorkonfiguration und ihrer physikalischen
Eigenschaften (Aufgaben~1 und~2), eine Bewertung frei verfügbarer
Open-Source-Sturzdatensätze (Aufgabe~3) sowie - als Hauptschwerpunkt - die
Entwicklung, das Training und die prototypische Integration eines spezifischen
Machine-Learning-Modells (Aufgabe~4).

\section{Umfang dieses Kapitels}
Dieses Dokument behandelt die Aufgaben~1 bis~3 des Projekts. Aufgabe~1
vergleicht die drei grundlegenden Sensorachsenkonfigurationen, Aufgabe~2
bewertet den Einfluss niedrigstufiger Sensoreigenschaften auf die
Erkennungsqualität und Aufgabe~3 sichtet die in Frage kommenden
Open-Source-Datensätze und begründet die Auswahl für die
Modellentwicklung in Aufgabe~4.

\section{Struktur}
Kapitel~\ref{ch:comparison} erläutert die physikalischen Grundlagen jeder
Sensorachsengruppe und leitet eine begründete Achsenempfehlung ab.
Kapitel~\ref{ch:properties} untersucht den Einfluss von fünf
Sensoreigenschaften auf die Erkennungsqualität.
Kapitel~\ref{ch:datasets} bewertet sieben Open-Source-Datensätze anhand
der in der Aufgabenstellung genannten Kriterien (Sensorposition,
Achsenkonfiguration und Abtastrate, Klassenabdeckung und Probanden\-diversität,
Verfügbarkeit) und begründet die Auswahl der Primär- und
Sekundärdatensätze für die Modellentwicklung.

%------------------------------------------------------------------
%------------------------------------------------------------------
\chapter{3-, 6- und 9-Achsen-Sensorsysteme zur Sturzerkennung}
\label{ch:comparison}

\section{Forschungsfrage und Vorgehensweise}
\label{sec:question}
Die zentrale Fragestellung dieses Kapitels lautet, ob eine höhere Anzahl
erfasster Achsen die Genauigkeit der Sturzerkennung statistisch signifikant
verbessert und bis zu welchem Punkt diese Verbesserung den zusätzlichen Aufwand
hinsichtlich Energieverbrauch, Rechenleistung und Komplexität rechtfertigt. Die
Aufgabenstellung verlangt ausdrücklich, dass die daraus abgeleitete Bewertung
durch Literatur oder Datenblätter gestützt wird und nicht lediglich auf
Behauptungen beruht. Die vorliegende Analyse kombiniert daher zwei voneinander
unabhängige Beweislinien: (i) begutachtete wissenschaftliche Studien, die
Sensitivität und Spezifität der jeweiligen Achsenkonfigurationen berichten,
sowie (ii) Herstellerdatenblätter repräsentativer Komponenten aus dem 6- und
9-Achsen-Bereich.

\section{Physikalische Bedeutung der Achsengruppen}
\label{sec:physics}
Die drei Achsenkonfigurationen unterscheiden sich hinsichtlich der
physikalischen Größen, die sie messen
(Tabelle~\ref{tab:axes-overview}):

\begin{table}[H]
\centering
\caption{Von 3-, 6- und 9-Achsen-Sensorkonfigurationen gemessene physikalische Größen.}
\label{tab:axes-overview}

\renewcommand{\arraystretch}{1.3}

\begin{tabularx}{\textwidth}{
>{\raggedright\arraybackslash}p{0.18\textwidth}
>{\raggedright\arraybackslash}p{0.27\textwidth}
>{\raggedright\arraybackslash}X}
\toprule
\textbf{Konfiguration} & \textbf{Sensoren} & \textbf{Gemessene Größe} \\
\midrule

3-Achsen
&
Dreiachsiger Beschleunigungssensor
&
Lineare Beschleunigung einschließlich der Erdbeschleunigung entlang der $x$-, $y$- und $z$-Achse, üblicherweise in g angegeben.
\\

6-Achsen (\ac{IMU})
&
Dreiachsiger Beschleunigungssensor + dreiachsiges Gyroskop
&
Lineare Beschleunigung zusammen mit der Winkelgeschwindigkeit (Roll-, Nick- und Gierrate), üblicherweise in $^\circ$/s angegeben.
\\

9-Achsen (\ac{MARG})
&
Dreiachsiger Beschleunigungssensor + dreiachsiges Gyroskop + dreiachsiges Magnetometer
&
Lineare Beschleunigung, Winkelgeschwindigkeit und der Magnetfeldvektor der Erde, wodurch mittels Sensorfusionsalgorithmen eine absolute Richtungsbestimmung und vollständige Orientierung ermöglicht werden.
\\

\bottomrule
\end{tabularx}
\end{table}

\section{3-Achsen-Systeme: Nur Beschleunigungssensor}
\label{sec:3axis}

\subsection{Erkennungsprinzip}
Der dominierende Ansatz bei 3-Achsen-Systemen berechnet den Betrag der
Beschleunigung (\ac{SVM}) und wendet Schwellwerte auf die charakteristische
Sturzsignatur an: eine Freifallphase, in der der Betrag gegen $0$\,g fällt,
gefolgt von einem ausgeprägten Aufprallmaximum. Ein häufig zitierter
Schwellwertalgorithmus von Bourke verwendet einen unteren Freifall-Schwellwert
von $0.65$\,g und einen oberen Aufprall-Schwellwert von $2.8$\,g, optional
ergänzt durch eine Überwachung der Körperhaltung nach dem Aufprall
(``Long-Lie'') \cite{bagala2012}.

\subsection{Berichtete Leistungsfähigkeit}
Schwellwertbasierte Algorithmen, die ausschließlich Beschleunigungssensoren
verwenden, erreichen unter Laborbedingungen typischerweise Sensitivitäts- und
Spezifitätswerte im Bereich von 80--95\,\%
\cite{dinh2021}. Bei sorgfältig simulierten Laborstürzen berichten einzelne
Algorithmen sogar nahezu perfekte Ergebnisse. Wurde derselbe
leistungsstärkste Algorithmus jedoch anhand einer Datenbank mit
\emph{realen} Stürzen älterer Menschen bewertet, sank die Sensitivität auf
$83\,\%$ (Spezifität $97\,\%$) \cite{bagala2012}. Diese Diskrepanz stellt
einen gut dokumentierten Realitätscheck für sämtliche Laborergebnisse in
diesem Forschungsgebiet dar.

\subsection{Zentrale Einschränkung: Falsch-positive Erkennungen}
Die ausschließliche Verwendung des Beschleunigungsbetrags erlaubt keine
zuverlässige Unterscheidung zwischen tatsächlichen Stürzen und sturzähnlichen
Aktivitäten des täglichen Lebens (\acp{ADL}), wie beispielsweise schnellem
Hinsetzen, Hinlegen, Laufen oder Springen, da all diese vergleichbar hohe
Beschleunigungsspitzen erzeugen können \cite{huynh2015}. Eine Absenkung des
Erkennungsschwellwertes, um mehr Stürze zu erfassen, erhöht unmittelbar die
Rate falsch-positiver Alarme, ein grundlegender Zielkonflikt zwischen
Sensitivität und Spezifität des Einzelmodalitätsansatzes
\cite{dinh2021}. Im Pflegekontext untergraben häufige Fehlalarme die
Akzeptanz der Nutzer und führen nachweislich zu Alarmmüdigkeit.

\section{6-Achsen-Systeme: Beschleunigungssensor und Gyroskop}
\label{sec:6axis}

\subsection{Welchen Beitrag das Gyroskop leistet}
Stürze gehen mit einer charakteristischen schnellen Änderung der
Körperorientierung (Rotation in Richtung der Horizontalen) einher, die bei
\acp{ADL} wie dem Hinsetzen nicht in derselben Kombination aus
Rotationsgeschwindigkeit und Beschleunigung auftritt. Die vom Gyroskop
gemessene Winkelgeschwindigkeit liefert daher ein unabhängiges
Unterscheidungsmerkmal genau dort, wo das Signal des Beschleunigungssensors
mehrdeutig ist.

\subsection{Berichtete Leistungssteigerungen}
Zwei unabhängige quantitative Ergebnisse kommen zu derselben Schlussfolgerung.
In einem auf der Brust getragenen System mit 702 Aufzeichnungen von
36 Probanden erhöhte die Ergänzung eines beschleunigungsbasierten Algorithmus
durch Gyroskop-Schwellwerte die \emph{Spezifität} von $82.72\,\%$ auf
$96.2\,\%$, während gleichzeitig eine Sensitivität von $96.3\,\%$ erhalten
blieb (Schwellwerte mittels ROC-Analyse optimiert)
\cite{huynh2015}. In einer Smartphone-basierten Studie mit frei variierenden
Gerätepositionen übertraf die Kombination aus Beschleunigungssensor und
Gyroskop die reine Beschleunigungssensor-Konfiguration um mehr als
$15$ Prozentpunkte hinsichtlich der Genauigkeit \cite{hakim2019}. Beide
Ergebnisse identifizieren die Reduktion falsch-positiver Erkennungen als den
wesentlichen Beitrag des Gyroskops.

\subsection{Grenzen des Nutzens des Gyroskops}
Der Zugewinn ist jedoch nicht uneingeschränkt. Studien, die die Mehrdeutigkeit
des Beschleunigungssensors bereits auf andere Weise kompensieren – etwa durch
zwei Sensorpositionen oder die Fusion mit einem Barometer-, berichten von
geringeren zusätzlichen Vorteilen durch das Gyroskop
\cite{arxiv2308_lstm,mdpi2025_realworld}. Der Vorteil einer
6-Achsen-Konfiguration ist daher insbesondere im Szenario mit nur einem Sensor
an einer einzigen Position am größten, genau der Einlagenkonfiguration, die
dieses Projekt verfolgt.

\section{9-Achsen-Systeme: Zusätzliches Magnetometer}
\label{sec:9axis}

\subsection{Funktion des Magnetometers}
Die auf dem Gyroskop basierende Orientierung wird durch numerische Integration
der Winkel-geschwindigkeit bestimmt, wodurch sich Messfehler im Laufe der Zeit
aufsummieren. Roetenberg et al.\ berichten, dass eine ausschließlich auf dem
Gyroskop basierende Integration innerhalb von $45$~s um etwa
$10^{\circ}$ abdriften kann \cite{roetenberg2005}. Das Magnetometer misst das
Magnetfeld der Erde und stellt dadurch eine \emph{absolute}
Richtungsreferenz bereit. In Kombination mit Beschleunigungssensor und
Gyroskop in einem Kalman- oder Madgwick-Filter korrigiert es diese Drift und
liefert eine stabile absolute Orientierung (Gier, Nick und Roll).

\subsection{Warum dies für die Sturzerkennung weniger relevant ist}
Ein Sturz ist ein kurzzeitiges Ereignis, die kritische Phase vom Verlust des
Gleichgewichts bis zum Aufprall dauert typischerweise nur ein bis wenige
Sekunden. Innerhalb dieses Zeitraums ist die Gyroskopdrift gering. Die
Driftkorrektur, welche den Hauptbeitrag des Magnetometers darstellt, kommt
daher in erster Linie der \emph{langfristigen} Orientierungsbestimmung zugute
und nicht der Erkennung eines kurzen, hochdynamischen Ereignisses. Die mit
sechs Achsen messbare relative Orientierungsänderung reicht aus, um die
Kinematik eines Sturzes zu charakterisieren.

\subsection{Nachteile: Magnetische Störungen in Innenräumen}
Messwerte von Magnetometern werden in Innenräumen systematisch verfälscht:
Ferromagnetische Materialien in Gebäudestrukturen, Haushaltsgeräten und
Möbeln verändern das lokale Magnetfeld, verursachen Richtungsfehler, die
häufig $5^{\circ}$ überschreiten, und erschweren dadurch mobile Anwendungen
\cite{pmc11991382,yang2024neurit,pubmed25347584}. Da ältere Menschen den
größten Teil ihrer Zeit in Innenräumen verbringen, erfordert ein robuster
Betrieb Kalibrierungs- und Störkompensationsalgorithmen, also zusätzlichen
Rechenaufwand \cite{pmc11991382,pubmed25347584} zusätzlich zum höheren
Energieverbrauch und Integrationsaufwand einer dritten Sensormodalität sowie
einer vollständigen 9-\ac{DoF}-Sensorfusion.

\section{Beispielkomponenten (Nachweis durch Datenblätter)}
\label{sec:datasheets}
Die drei in der Aufgabenstellung genannten repräsentativen Komponenten werden
in Tabelle~\ref{tab:components} verglichen. Sämtliche Angaben sind typische
Werte unter den jeweiligen Herstellerbedingungen und müssen für den finalen
Bericht anhand der Originaldatenblätter erneut überprüft werden.

\begin{table}[H]
\centering
\caption{Repräsentative Komponenten auf 6- und 9-Achsen-Ebene.}
\label{tab:components}

\footnotesize
\setlength{\tabcolsep}{4pt}
\renewcommand{\arraystretch}{0.95}

\begin{tabularx}{\textwidth}{
>{\raggedright\arraybackslash}p{2.1cm}
>{\centering\arraybackslash}p{1.6cm}
>{\raggedright\arraybackslash}X
>{\centering\arraybackslash}p{2cm}
>{\raggedright\arraybackslash}X
}
\toprule
\textbf{Komponente} &
\textbf{Konfig.} &
\textbf{Messbereiche} &
\textbf{Stromaufnahme} &
\textbf{Besonderheiten} \\
\midrule

MPU-6050 \cite{invensense_mpu6050} &
6-Achsen \ac{IMU} &
$\pm2/4/8/16$\,g; $\pm250$--$2000^{\circ}$/s &
$\approx3.9$\,mA &
Klassische kostengünstige Referenz mit 16-Bit-\acp{ADC}; alle Sensoren aktiv; Referenz für den 6-Achsen-Vergleich. \\

LSM6DSOX \cite{st_lsm6dsox} &
6-Achsen \ac{IMU} &
$\pm2$--$16$\,g; $\pm125$--$2000^{\circ}$/s &
0.55\,mA &
Hochleistungsmodus; etwa $7\times$ geringere Stromaufnahme als der MPU-6050; integrierter \ac{MLC} und hardwarebasierte Freifallerkennung. \\

BHI360 \cite{bosch_bhi360} &
6-Achsen \ac{IMU}\\
(9-\ac{DoF} externes Magnetometer) &
Geräteabhängig &
$<600$\,\textmu A &
Führt die BSX-Fusionsbibliothek direkt auf dem Chip aus und entlastet dadurch den Hauptprozessor von der Sensorfusion. \\

\bottomrule
\end{tabularx}
\end{table}

Der Effizienzunterschied zwischen dem MPU-6050 (Generation um 2010) und dem
LSM6DSOX/BHI360 zeigt außerdem, dass die \emph{Sensorgeneration} den
Energiebedarf ebenso stark beeinflussen kann wie die Anzahl der Achsen.

\section{Vergleich auf einen Blick}
\label{sec:glance}

\begin{table}[H]
\centering
\caption{Direkter Vergleich der Achsenkonfigurationen für die Sturzerkennung.}
\label{tab:glance}
\small
\begin{tabularx}{\textwidth}{l X X X}
\toprule
 & \textbf{3-Achsen} & \textbf{6-Achsen} & \textbf{9-Achsen} \\
\midrule
Erfasst
& Aufprallmaximum, Freifallphase (\ac{SVM})
& + Rotationsgeschwindigkeit $\rightarrow$ Sturzkinematik
& + absolute Orientierung / Himmelsrichtung \\
Stärken
& Minimaler Aufwand hinsichtlich Kosten, Energie und Komplexität
& Großer Spezifitätsgewinn ($82.7\rightarrow96.2\,\%$ \cite{huynh2015}; über $15$ Prozentpunkte Genauigkeit \cite{hakim2019})
& Driftkorrigierte langfristige Orientierung \cite{roetenberg2005} \\
Schwächen
& Falsch-positive Erkennungen durch sturzähnliche \acp{ADL}
\cite{dinh2021,huynh2015}; geringere Sensitivität bei realen Stürzen
\cite{bagala2012}
& Höhere Stromaufnahme; Drift spielt nur deshalb kaum eine Rolle, weil Stürze kurz sind
& Magnetische Störungen in Innenräumen
\cite{pmc11991382,yang2024neurit,pubmed25347584}; zusätzlicher Energie-, Rechen- und Kalibrierungsaufwand \\
Nutzen für die Sturzerkennung
& Basislösung, praktikabel, aber anfällig für Alarmmüdigkeit
& Hoch, behebt den wesentlichen Schwachpunkt
& Gering, korrigiert einen Fehler, der sich innerhalb von $1$-$2$\,s kaum aufbaut \\
\bottomrule
\end{tabularx}
\end{table}

\section{Begründetes Fazit}
\label{sec:conclusion}

\textbf{Von drei auf sechs Achsen: eindeutig gerechtfertigt.} Der Übergang
adressiert den dokumentierten Hauptschwachpunkt der ausschließlich auf
Beschleunigungssensoren basierenden Sturzerkennung - falsch-positive
Erkennungen durch sturzähnliche \acp{ADL} - mit gemessenen
Spezifitätssteigerungen von etwa $\approx83\,\%$ auf
$\approx96\,\%$ \cite{huynh2015} sowie Genauigkeitssteigerungen von mehr als
15 Prozentpunkten \cite{hakim2019}. Dies geschieht zum Preis einer einzigen
zusätzlichen \ac{MEMS}-Sensormodalität, deren Stromaufnahme bei modernen
Bauteilen gering ausfällt ($0.55$\,mA für eine vollständige
6-Achsen-LSM6DSOX \cite{st_lsm6dsox}).

\textbf{Von sechs auf neun Achsen: für diesen Anwendungsfall nicht
verhältnismäßig.} Der Hauptbeitrag des Magnetometers, absolute
Richtungsbestimmung und langfristige Driftkorrektur
\cite{roetenberg2005} adressiert ein Problem, das bei einem
$1$-$2$ Sekunden dauernden Sturzereignis praktisch nicht relevant ist,
führt jedoch gleichzeitig ein neues Problem ein: Richtungsfehler durch
ferromagnetische Verzerrungen genau in den Innenraumumgebungen, in denen
ältere Menschen leben
\cite{pmc11991382,yang2024neurit,pubmed25347584}, sowie zusätzlichen
Energieverbrauch, Rechenaufwand und Kalibrierungsbedarf. Keine der in dieser
Literaturübersicht betrachteten Studien zeigt einen Genauigkeitsgewinn durch
das Magnetometer, der mit dem Beitrag des Gyroskops vergleichbar wäre.

\textbf{Empfehlung:} Eine 6-Achsen-Konfiguration stellt den besten Kompromiss
für die einlagenbasierte Sturzerkennung dar. Dies wird zusätzlich dadurch
unterstützt, dass die führenden öffentlich verfügbaren Sturzdatensätze
(SisFall, KFall) selbst auf 6-Achsen-Aufzeichnungen basieren, wodurch zugleich
die Kompatibilität der Trainingsdaten für das im Rahmen von Aufgabe~4 zu
entwickelnde Machine-Learning-Modell sichergestellt wird.


%------------------------------------------------------------------
\chapter{Einfluss der Sensoreigenschaften auf die Erkennungsqualität}
\label{ch:properties}

In Kapitel~\ref{ch:comparison} wurde die Wahl einer sechsachsigen \ac{IMU} auf
der Ebene begründet, \emph{welche} Achsen erfasst werden sollen. Das
vorliegende Kapitel ergänzt diese Argumentation um eine zweite Ebene und
fragt, \emph{wie präzise} diese Achsen erfasst werden müssen. Betrachtet
werden fünf niedrigstufige Datenblattgrößen, nämlich Messbereich,
Abtastrate, Bit-Auflösung, Rauschdichte und Drift, sowie deren jeweiliger
Einfluss auf die Sturzerkennung. Gemäß der Aufgabenstellung wird der
Umfang bewusst kompakt gehalten und stützt sich vollständig auf die
Fachliteratur; eigene Messreihen sind nicht vorgesehen.

\section{Messbereich und Signal-Clipping}
\label{sec:range}
Der Messbereich legt fest, wie große Beschleunigungen ein Sensor darstellen
kann, bevor er in die Sättigung gerät. Ein Sturz besteht aus zwei
dynamischen Phasen mit stark unterschiedlicher Amplitude. In der freien
Fallphase sinkt der Betrag der gemessenen Beschleunigung gegen $0$\,g,
worauf innerhalb weniger Millisekunden der Aufprall folgt, dessen Spitzenwert
typischerweise ein Mehrfaches der Erdbeschleunigung erreicht. In der
Literatur häufig zitierte Schwellwertverfahren setzen den oberen
Impact-Schwellwert an der Hüfte im Bereich von $2{,}0$\,g bis $2{,}8$\,g
an \cite{bagala2012,us10152869}. An peripheren Körperstellen fallen die
lokalen Spitzenwerte allerdings deutlich höher aus; Untersuchungen zur
Tibia-Beschleunigung im Alltag berichten Werte, die im Knöchelbereich
$16$\,g bis $24$\,g überschreiten können \cite{cornelissen2023clipping}.

Wird der Messbereich zu klein gewählt, kommt es zum sogenannten Clipping.
Das Analogsignal wird auf den größten darstellbaren Wert abgeschnitten,
und in der Zeitreihe zeigen sich charakteristische flache Plateaus. Dabei
wird gerade dasjenige Merkmal verzerrt, das für die Sturzerkennung die
höchste Trennschärfe aufweist, nämlich der Impact-Peak.
Vergleichende Messungen belegen einen Genauigkeitsverlust bei der
Peak-Tibial-Acceleration von etwa $6{,}25\,\%$ bei $\pm 16$\,g und
$28{,}17\,\%$ bei $\pm 8$\,g gegenüber einer Referenz mit $\pm 70$\,g
\cite{cornelissen2023clipping}. Ein Messbereich von $\pm 2$\,g, wie er in
einfachen Konsumbausteinen üblich ist, reicht daher an keiner
Körperposition aus. $\pm 8$\,g stellt eine praxistaugliche Untergrenze
dar, während $\pm 16$\,g, wie im SisFall-Datensatz verwendet
\cite{sucerquia2017sisfall}, insbesondere für periphere Positionen wie
die geplante Anbringung in der Schuhsohle einen sinnvollen
Sicherheitspuffer bietet.

\section{Abtastrate}
\label{sec:samplingrate}
Die Abtastrate bestimmt, mit welcher zeitlichen Auflösung sich die
charakteristische Signatur eines Sturzes rekonstruieren lässt.
Menschliche Bewegungen sind spektral von tiefen Frequenzen dominiert, und
die kritische Phase eines Sturzes umfasst nur einige hundert Millisekunden.
Nach dem Nyquist-Kriterium sind daher bereits Abtastraten im Bereich
einiger zehn Hertz ausreichend, um das relevante Frequenzband abzudecken.

Zwei unabhängig voneinander durchgeführte empirische Arbeiten kommen zu
einem übereinstimmenden Ergebnis. Villar et al.\ zeigen, dass eine
Abtastfrequenz von $15$\,Hz bis $20$\,Hz genügt, um mit einem
Convolutional Neural Network eine Sensitivität und Spezifität von über
$95\,\%$ zu erreichen, während höhere Frequenzen keinen nennenswerten
Zusatzgewinn liefern \cite{villar2022sampling}. Eine aktuellere
Untersuchung auf dem SisFall-Datensatz trainierte fünf Klassifikatoren
bei $10$, $20$, $50$ und $100$\,Hz. Als bester Kompromiss wurde ein
CNN-LSTM-Modell bei $20$\,Hz identifiziert, das eine Genauigkeit von
$98{,}9\,\%$, eine Sensitivität von $96{,}7\,\%$ und eine
Spezifität von $99{,}6\,\%$ erreicht und dabei gleichzeitig das
übertragene Datenvolumen deutlich reduziert \cite{diaz2026sampling}.
Werte unterhalb von $10$\,Hz beginnen, kritische Details der
Aufprallphase zu verlieren, während Werte oberhalb von etwa $100$\,Hz
im Wesentlichen nur die Datenrate und den Energiebedarf erhöhen, ohne
dass die Klassifikationsleistung weiter zunimmt.

Für ein sohlenbasiertes System, das ein batteriebetriebenes Wearable mit
einem auf begrenzter Hardware ausgeführten ML-Modell kombiniert, stellt
somit eine Abtastrate zwischen $20$\,Hz und $50$\,Hz einen aus praktischer
Sicht sinnvollen Arbeitspunkt dar.

\section{Bit-Auflösung}
\label{sec:resolution}
Digital ausgegebene \ac{MEMS}-Beschleunigungssensoren quantisieren das
Analogsignal mit einer festen Anzahl von Bits, in der Regel zwischen zwölf
und sechzehn. Für die nominelle Auflösung $q$ in Einheiten von g gilt
$q = \mathrm{FSR}/2^{N}$, wobei $\mathrm{FSR}$ die Vollbereichsspanne und
$N$ die Bit-Breite bezeichnet. Für einen 16-Bit-\ac{ADC} bei $\pm 16$\,g
folgt daraus $q \approx 0{,}5$\,mg pro Code. Eine höhere Auflösung
erlaubt eine feinere Unterscheidung geringfügiger Bewegungen, etwa
langsames Hinsetzen im Vergleich zum Einleiten eines Sturzes, was der
merkmalsbasierten Machine-Learning-Stufe in Aufgabe~4 zugutekommt.

Die nominelle Auflösung stellt allerdings eine obere Schranke dar. Die
effektive Bitanzahl richtet sich nach dem Verhältnis von
Vollbereichsspanne zum integrierten Sensorrauschen. Veröffentlichte
Beispiele belegen, dass ein als 16-Bit deklarierter Baustein nach Abzug
des sensoreigenen Rauschens nur etwa $8{,}5$ rauschfreie Bits liefert
\cite{nxp2013an4075}. In der Praxis gehören 12- bis 16-Bit-Wandler in
Wearables zur Sturzerkennung zum Stand der Technik und stellen keinen
limitierenden Faktor der Klassifikationsgüte dar. Praktisch relevanter
ist die im folgenden Abschnitt behandelte Rauschdichte.

\section{Rauschdichte}
\label{sec:noise}
Die Rauschdichte $n_{d}$, in Datenblättern typischerweise in
$\mu\mathrm{g}/\sqrt{\mathrm{Hz}}$ für den Beschleunigungssensor
beziehungsweise in $\mathrm{mdps}/\sqrt{\mathrm{Hz}}$ für das Gyroskop
angegeben, quantifiziert das dem Nutzsignal überlagerte stochastische
Grundrauschen. Für das effektive RMS-Rauschen am Klassifikatoreingang
gilt

\begin{equation}
\sigma_{\mathrm{RMS}} \;=\; n_{d} \cdot \sqrt{\mathrm{BW}_{\mathrm{NEB}}},
\label{eq:noise}
\end{equation}
wobei $\mathrm{BW}_{\mathrm{NEB}}$ die rauschäquivalente Bandbreite des
digitalen Filters bezeichnet \cite{arar2022noise}. Für einen typischen
Konsumbaustein mit $n_{d} \approx 126\,\mu\mathrm{g}/\sqrt{\mathrm{Hz}}$
und einer Bandbreite von $200$\,Hz ergibt sich hieraus
$\sigma_{\mathrm{RMS}} \approx 1{,}8$\,mg \cite{arar2022noise}. Eine
verbreitete Faustregel fordert, dass das kleinste zu detektierende
Merkmal mindestens eine Größenordnung über $\sigma_{\mathrm{RMS}}$
liegen sollte \cite{arar2022noise}.

Diese Anforderung ist für die Sturzerkennung leicht zu erfüllen. Die
dominanten Impact-Peaks liegen im Bereich mehrerer g, sodass das
Grundrauschen die Aufprallphase in der Praxis nicht beeinflusst. Relevant
wird die Rauschdichte hingegen in der Merkmalsextraktion einer
Machine-Learning-Pipeline (vgl.\ Aufgabe~4), in der feine Statistiken wie
Varianz, Kurtosis oder Ableitungen der Signalmagnitude herangezogen
werden, um sturzähnliche Alltagsaktivitäten von tatsächlichen
Stürzen zu trennen. Rauscharme Bausteine wie der LSM6DSOX
\cite{st_lsm6dsox} bringen an dieser Stelle einen kleinen, aber
reproduzierbaren Vorteil für den Klassifikator, und dies ohne
relevanten Zusatzaufwand, da sie zugleich einen geringeren
Stromverbrauch aufweisen.

\section{Drift, insbesondere beim Gyroskop}
\label{sec:drift}
Das Ausgangssignal eines Gyroskops enthält einen langsam
veränderlichen Bias, dessen spektrale Fluktuationen im tiefen
Frequenzbereich von einem Anteil mit $1/f$-Charakteristik dominiert werden
\cite{vagner2013allan}. Die zugehörige Kenngröße, die
Bias-Instabilität, wird üblicherweise über die Allan-Varianz
quantifiziert. Für typische kostengünstige \ac{MEMS}-Gyroskope liegt
sie in der Größenordnung von $3{,}6\,^{\circ}/\mathrm{h}$
\cite{gakstatter2017biasdrift}. Bausteine der MPU-6050-Klasse zeigen auf
allen drei Achsen ein vergleichbares Verhalten
\cite{ilyas2020mpu6050allan}.

Ob dieser Drift praxisrelevant wird, hängt entscheidend vom
Beobachtungsfenster ab. Bei einer Integration der Winkelgeschwindigkeit zur
Bestimmung der Orientierung akkumuliert sich der Fehler mit der Zeit.
Roetenberg et al.\ weisen für eine rein gyroskopische Integration einen
bias-bedingten Orientierungsfehler von rund $10^{\circ}$ innerhalb von
$45$~Sekunden nach \cite{roetenberg2005}. Für ein Sturzereignis, dessen
relevante Zeitspanne im Bereich von ein bis zwei Sekunden liegt, ist der
Driftbeitrag jedoch verschwindend gering im Vergleich zu den während
des Sturzes auftretenden Rotationsraten. Dieser Befund deckt sich mit der
Schlussfolgerung aus Kapitel~\ref{ch:comparison}, wonach die
driftkompensierende Funktion des Magnetometers für die eigentliche
Sturzerkennung nur einen marginalen Nutzen liefert.

\section{Zusammenfassung}
\label{sec:synthesis}
Die fünf betrachteten Eigenschaften wirken nicht isoliert, sondern greifen
ineinander (Tabelle~\ref{tab:properties}). Für einen Machine-Learning-
basierten Sturzdetektor in Sohlenkonfiguration ergibt sich hieraus als
empfohlener Arbeitspunkt eine sechsachsige \ac{IMU} mit einem
Beschleunigungs-Messbereich von $\pm 16$\,g, einem Gyroskop-Messbereich
von $\pm 1000$ bis $\pm 2000\,^{\circ}/\mathrm{s}$, einer Abtastrate
zwischen $20$\,Hz und $50$\,Hz, mindestens $12$~Bit Auflösung, geringer
Rauschdichte (unterhalb von etwa
$100\,\mu\mathrm{g}/\sqrt{\mathrm{Hz}}$) sowie einer nicht weiter
kompensierten Gyroskop-Drift, da deren Beitrag über die kurze
Sturzdauer vernachlässigbar bleibt. Diese Anforderungen werden vom in
Kapitel~\ref{ch:comparison} diskutierten LSM6DSOX abgedeckt und stehen im
Einklang mit der Aufnahmecharakteristik des SisFall-Datensatzes, der in
den Aufgaben~3 und~4 verwendet wird.

\begin{table}[H]
\centering
\caption{Einfluss der Sensoreigenschaften auf die Sturzerkennungsqualität.}
\label{tab:properties}
\small
\begin{tabularx}{\textwidth}{l X X}
\toprule
\textbf{Eigenschaft} & \textbf{Primärer Effekt} & \textbf{Empfohlener Bereich} \\
\midrule
Messbereich       & Verhindert Clipping des Impact-Peaks; ein zu kleiner Bereich verzerrt das trennschärfste Merkmal \cite{cornelissen2023clipping,us10152869} & mindestens $\pm 8$\,g, empfohlen $\pm 16$\,g \\
Abtastrate        & Zeitliche Rekonstruktion der Sturzsignatur; höhere Raten erhöhen den Energiebedarf ohne Genauigkeitsgewinn \cite{villar2022sampling,diaz2026sampling} & $20$\,Hz bis $50$\,Hz \\
Bit-Auflösung   & Feinauflösung der Merkmalsextraktion; effektive Bitzahl durch das Rauschen begrenzt \cite{nxp2013an4075} & mindestens $12$~Bit \\
Rauschdichte      & Bestimmt die minimal auflösbare Signalgröße; wirkt sich vor allem auf die Unterscheidung von ADL und Präaufprall aus \cite{arar2022noise} & $\leq 100\,\mu\mathrm{g}/\sqrt{\mathrm{Hz}}$ \\
Gyroskop-Drift    & Langzeitfehler der Orientierungsbestimmung; auf der Zeitskala eines Sturzes vernachlässigbar \cite{vagner2013allan,roetenberg2005} & für kurze Ereignisse unkritisch \\
\bottomrule
\end{tabularx}
\end{table}

%------------------------------------------------------------------
\chapter{Ausblick auf Datensätze für das ML-Modell}
\label{ch:datasets}

Für die Entwicklung des eigenen Machine-Learning-Modells in Aufgabe~4 stehen
mehrere frei verfügbare Sturzdatensätze zur Auswahl. Die Aufgabenstellung
nennt mindestens sieben Kandidaten, die im Folgenden hinsichtlich ihrer
Eignung für den angestrebten Anwendungsfall bewertet werden: eine
schuhsohlenintegrierte (\emph{Insole}-)Sturzerkennung bei älteren,
gegebenenfalls kognitiv eingeschränkten Personen. Die Bewertungskriterien
orientieren sich an der Aufgabenstellung und umfassen die Sensorposition am
Körper und deren Übertragbarkeit auf eine Fußsohle, die
Achsenkonfiguration und Abtastrate im Vergleich zur geplanten Sensorik, die
Abdeckung verschiedener Sturz- und ADL-Klassen, die Diversität der Probanden
sowie die Verfügbarkeit und Zugänglichkeit der Rohdaten.

Die Aufgabenstellung enthält bereits eine erste Einschätzung zur Eignung der
sieben Datensätze. Diese wurde im Rahmen der vorliegenden Bearbeitung anhand
der Originalpublikationen und der offiziellen Datensatzquellen eigenständig
geprüft, vertieft und dort präzisiert oder korrigiert, wo die tatsächlichen
Datenblätter von der Aufgabenstellung abweichen. Die entsprechenden
Anpassungen sind in den nachfolgenden Abschnitten kenntlich gemacht.

\section{Übersicht der Datensätze}
\label{sec:dataset_overview}

Tabelle~\ref{tab:datasets} fasst die wesentlichen Eckdaten der sieben
Datensätze zusammen. Im Anschluss wird jeder Datensatz einzeln hinsichtlich
seiner Stärken und Einschränkungen für den vorliegenden Anwendungsfall
eingeordnet.

\begin{table}[H]
\centering
\caption{Vergleich der sieben zu bewertenden Sturzdatensätze.}
\label{tab:datasets}
\footnotesize
\setlength{\tabcolsep}{3pt}
\renewcommand{\arraystretch}{1.15}
\begin{tabularx}{\textwidth}{
>{\raggedright\arraybackslash}p{1.6cm}
>{\raggedright\arraybackslash}p{2.3cm}
>{\centering\arraybackslash}p{1.2cm}
>{\centering\arraybackslash}p{1.0cm}
>{\centering\arraybackslash}p{1.4cm}
>{\raggedright\arraybackslash}X
}
\toprule
\textbf{Datensatz} &
\textbf{Sensorik / Position} &
\textbf{Achsen} &
\textbf{Hz} &
\textbf{Probanden} &
\textbf{Eignung für Insole} \\
\midrule

SisFall \cite{sucerquia2017sisfall}
& 2 Beschl.-Sensoren (ADXL345 $\pm 16$\,g, MMA8451Q $\pm 8$\,g) + Gyroskop (ITG3200 $\pm 2000\,^{\circ}/\mathrm{s}$), am Gürtel (Hüfte)
& 6 (2$\times$3 + 3)
& 200
& 38 (23 jung, 15 ältere)
& Gut geeignet: passende Achsenkonfiguration, hohe Abtastrate, etablierter Benchmark mit vielen Vergleichswerten in der Literatur. \\

MobiFall / MobiAct \cite{mobifall2014,mobiact2016}
& Smartphone (Samsung Galaxy S3) in der Hosentasche
& 6--9 (variabel)
& $\approx$\,20--87
& 11 / 57
& Bedingt geeignet: Smartphone-Sensorik (oft nur $\pm 2$\,g) unterscheidet sich von einer Insole; Trageposition und Messbereich weichen ab. \\

UMAFall \cite{umafall2017}
& 4$\times$ SensorTag (MPU-9250) an Brust, Hüfte, Handgelenk, Knöchel + Smartphone in der Tasche
& 9
& 20 (Tags) / 200 (Smartphone)
& 19 (19--68~J.)
& Bedingt geeignet: Mehrknoten-Aufbau interessant, aber IMU-Abtastrate nur 20\,Hz; Knöcheldaten vorhanden, jedoch nur 3 Sturztypen. \\

KFall \cite{kfall2021}
& 9-Achsen-IMU (MTw~Awinda) am unteren Rücken; bereitgestellt werden Beschl., Winkelgeschw.\ und Euler-Winkel
& 6 + Euler
& 100
& 32
& Gut geeignet: hohe Abtastrate, 15 Sturztypen, zeitliche Annotation des Aufprallzeitpunkts (Pre-Impact), umfangreicher Benchmark. \\

FallAllD \cite{fallalld2020}
& 3$\times$ IMU (LSM9DS1) an Nacken, Handgelenk und Taille; Beschl.\ $\pm 8$\,g, Gyro $\pm 2000\,^{\circ}/\mathrm{s}$, Mag, Baro
& 9 (+ Baro)
& 238 / 476 (Beschl.)
& 15
& Bedingt geeignet: Taillendaten nutzbar, hohe Abtastrate, aber nur 15 (junge) Probanden und Messbereich $\pm 8$\,g kann an peripheren Positionen clippen (vgl.\ Abschnitt~\ref{sec:range}). \\

UniMiB SHAR \cite{unimib2017}
& Smartphone (Samsung Galaxy Nexus) in der Hosentasche
& 3
& 50
& 30 (18--60~J.)
& Eingeschränkt geeignet: rein 3-achsig (nur Beschleunigungssensor), daher für einen theoretischen 3- vs.\ 6-Achsen-Vergleich hilfreich, aber nicht für ein realistisches 6-Achsen-Modell. \\

WEDA-FALL \cite{wedafall2023}
& Fitbit Sense Smartwatch am Handgelenk
& 6
& 50
& 25 (14 jung, 11 ältere $>$80\,J.)
& Bedingt geeignet: Handgelenk-Position weicht stark von der Fußposition ab; niedrige Abtastrate (50\,Hz); enthält jedoch Daten älterer Probanden, was für Kreuzvalidierung nützlich sein kann. \\

\bottomrule
\end{tabularx}
\end{table}

\section{Einzelbewertung}
\label{sec:dataset_details}

\textbf{SisFall} wurde an der Universidad de Antioquia erhoben und umfasst
4505 Aufzeichnungen (2707 ADL, 1798 simulierte Stürze) von 38 Probanden,
darunter 15 Personen im Alter von 60 bis 75 Jahren
\cite{sucerquia2017sisfall}. Die Sensoreinheit wurde am Gürtel (Hüfte)
befestigt und besteht aus \emph{zwei} Beschleunigungssensoren (ADXL345 mit
$\pm 16$\,g und MMA8451Q mit $\pm 8$\,g) sowie einem Gyroskop (ITG3200
mit $\pm 2000\,^{\circ}/\mathrm{s}$), aufgezeichnet mit einer Abtastrate
von 200\,Hz. Die redundante Beschleunigungsmessung entspricht der Angabe
der Aufgabenstellung von ,,zwei Beschleunigungssensoren + Gyroskop`` und
wurde in einer früheren Fassung dieser Arbeit unzutreffend als
Einzelsensor beschrieben. SisFall ist der am häufigsten zitierte Benchmark
für die sensorbasierte Sturzerkennung, was eine direkte Vergleichbarkeit
der eigenen Ergebnisse mit der Literatur ermöglicht. Die
6-Achsen-Konfiguration (Beschleunigung + Gyroskop) stimmt mit der in
Kapitel~\ref{ch:comparison} empfohlenen Achsenwahl überein. Einschränkend
ist anzumerken, dass die älteren Probanden nur ADL durchführten und keine
simulierten Stürze, sodass die Sturzdaten ausschließlich von jungen
Erwachsenen stammen.

\textbf{MobiFall und MobiAct} basieren auf Smartphone-Sensorik und wurden an
der Technischen Hochschule Kreta erhoben \cite{mobifall2014,mobiact2016}.
MobiAct umfasst 57 Probanden, 4 Sturztypen und 9 ADL; die Daten wurden mit
dem eingebauten Beschleunigungssensor und Gyroskop eines Samsung Galaxy S3
aufgezeichnet. Der Messbereich der Smartphone-Beschleunigungssensoren liegt
häufig nur bei $\pm 2$\,g, was für die Sturzerkennung zu Clipping führen
kann (vgl.\ Abschnitt~\ref{sec:range}). Darüber hinaus unterscheidet sich
die Trageposition (Hosentasche) grundlegend von einer Insole, sodass eine
Übertragbarkeit der Bewegungsmuster auf die Fußsohle nicht ohne Weiteres
gegeben ist.

\textbf{UMAFall} setzt einen Mehrknoten-Aufbau mit vier SensorTags
(MPU-9250, je Beschleunigungssensor, Gyroskop und Magnetometer) und einem
Smartphone ein \cite{umafall2017}. Durch die gleichzeitige Erfassung an Brust,
Hüfte, Handgelenk und Knöchel ermöglicht dieser Datensatz die
Untersuchung positionsabhängiger Effekte. Die Abtastrate der SensorTags
beträgt allerdings nur 20\,Hz, und es sind lediglich drei Sturztypen
enthalten. Für den vorliegenden Anwendungsfall bieten die Knöcheldaten
zwar eine gewisse Nähe zur Fußposition, die geringe Abtastrate und
begrenzte Sturzdiversität schränken den praktischen Nutzen jedoch ein.

\textbf{KFall} wurde am KAIST in Südkorea erhoben und umfasst 5075
Aufzeichnungen (2729 ADL, 2346 Stürze) von 32 Probanden
\cite{kfall2021}. Zum Einsatz kam eine 9-Achsen-IMU (Xsens MTw~Awinda) am
unteren Rücken; die veröffentlichten Datenfiles enthalten die
3-Achsen-Beschleunigung, die 3-Achsen-Winkelgeschwindigkeit sowie die per
Sensorfusion abgeleiteten Euler-Winkel bei 100\,Hz. Für den Vergleich
mit einer geplanten reinen 6-Achsen-Sensorik können die Euler-Kanäle
verworfen und lediglich Beschleunigung und Winkelgeschwindigkeit
verwendet werden. Eine Besonderheit von KFall ist die zeitliche
Annotation des Sturzbeginns und des Aufprallzeitpunkts anhand
synchronisierter Videoaufnahmen, was den Datensatz insbesondere für die
Pre-Impact-Erkennung wertvoll macht. KFall bietet damit eine gute
Ergänzung zu SisFall, da die simulierten Sturztypen aus SisFall
übernommen wurden, die Abtastrate hinreichend hoch ist und die
Pre-Impact-Annotation Zusatzinformationen für die Merkmalsextraktion
liefert. Einschränkend ist, dass sämtliche Probanden junge Erwachsene
waren.

\textbf{FallAllD} enthält Daten von drei gleichzeitig getragenen IMUs
(LSM9DS1), die am Nacken, am Handgelenk und an der Taille befestigt waren
\cite{fallalld2020}. Damit weicht die Sensorposition von der Angabe
in der Aufgabenstellung (Hals, Brust, Taille) ab; entscheidend ist die Position
am Handgelenk anstelle der Brust, wie die offizielle Beschreibung des
Datensatzes auf IEEE~DataPort bestätigt. Die Abtastrate ist mit 238\,Hz
für Beschleunigungssensor und Gyroskop hoch, die Sensoren umfassen zusätzlich
ein Magnetometer und ein Barometer. Der Datensatz deckt zahlreiche Sturz-
und ADL-Typen ab und ist hinsichtlich der Bewegungsvielfalt umfangreicher
als die meisten anderen betrachteten Kandidaten. Einschränkend sind der
gegenüber SisFall geringere Messbereich von $\pm 8$\,g sowie die kleine
Probandenzahl von 15 (jungen) Personen. Für den vorliegenden Anwendungsfall
sind vor allem die Aufzeichnungen am Taillensensor relevant, da sie in der
Kinematik am ehesten den bisherigen Benchmark-Datensätzen entsprechen.

\textbf{UniMiB SHAR} ist ein reiner 3-Achsen-Datensatz
(Beschleunigungssensor, 50\,Hz) von Smartphone-Aufzeichnungen aus der
Hosentasche \cite{unimib2017}. Er umfasst 11\,771 Samples von 30 Probanden mit
9 ADL und 8 Sturztypen. Da kein Gyroskop vorhanden ist, eignet sich dieser
Datensatz nicht für die Entwicklung eines 6-Achsen-Modells, kann aber
hilfsweise als Vergleichsbasis für die Frage herangezogen werden, welchen
Mehrwert die sechste Achse bringt (vgl.\ Kapitel~\ref{ch:comparison}).

\textbf{WEDA-FALL} wurde mit einer Fitbit Sense Smartwatch am Handgelenk
aufgenommen und enthält Beschleunigungs- und Gyroskop-Daten mit einer
nativen Abtastrate von 50\,Hz \cite{wedafall2023}. Nach Angabe der
Aufgabenstellung stellt WEDA-FALL zusätzlich vordefinierte
Downsampling-Varianten bei 40, 25, 10 und 5\,Hz bereit. Diese
Mehrfach-Abtastraten bilden ein zusätzliches Argument für die Aufnahme
von WEDA-FALL in die engere Auswahl: Der Einfluss der in
Kapitel~\ref{ch:properties} diskutierten Abtastrate lässt sich damit an
identischen Bewegungssequenzen kontrolliert isolieren, was in keinem der
übrigen Datensätze möglich ist. Eine weitere Stärke ist die Einbeziehung
von 11 älteren Probanden (über 80 Jahre), was unter den hier betrachteten
Datensätzen einzigartig ist. Die Handgelenk-Position weicht jedoch stark
von einer Fußsohle ab, und die vergleichsweise niedrige native
Abtastrate begrenzt die Auflösung hochdynamischer Phasen. WEDA-FALL ist
daher vor allem für handgelenkbasierte Szenarien relevant und bietet
sich für eine Kreuzvalidierung des Modells sowie für gezielte
Ablationsstudien zur Abtastrate an.

\section{Verfügbarkeit und Zugänglichkeit}
\label{sec:dataset_availability}

Neben den inhaltlichen Kriterien nennt die Aufgabenstellung ausdrücklich
die \emph{Verfügbarkeit und Zugänglichkeit} der Datensätze als
Bewertungskriterium. Tabelle~\ref{tab:availability} stellt für jeden der
sieben Kandidaten die offizielle Bezugsquelle, die Art des Zugangs sowie
die Lizenz zusammen. Der Zugang reicht dabei von einem unmittelbaren
Direktdownload ohne Registrierung (UMAFall, KFall, WEDA-FALL) über
einfache Kontoerstellungen (FallAllD via IEEE~DataPort) bis hin zu einem
formal zu unterzeichnenden Nutzungsvertrag mit dem herausgebenden Institut
(MobiFall/MobiAct). Für die zeitliche Planung des ML-Teils in Aufgabe~4
ist insbesondere zu berücksichtigen, dass die Freigabe der
MobiFall/MobiAct-Daten eine mehrtägige Bearbeitungszeit erfordert und
daher frühzeitig zu beantragen ist. Für SisFall wurde im Rahmen dieser
Arbeit festgestellt, dass die in mehreren Publikationen zitierte URL
zeitweise nicht erreichbar war; die Daten sind derzeit jedoch wieder über
die Projektseite des SISTEMIC-Labors an der Universidad de Antioquia
abrufbar.

\begin{table}[H]
\centering
\caption{Verfügbarkeit und Zugänglichkeit der betrachteten Datensätze.}
\label{tab:availability}
\footnotesize
\setlength{\tabcolsep}{3pt}
\renewcommand{\arraystretch}{1.2}
\begin{tabularx}{\textwidth}{
>{\raggedright\arraybackslash}p{1.7cm}
>{\raggedright\arraybackslash}X
>{\raggedright\arraybackslash}p{3.0cm}
>{\raggedright\arraybackslash}p{2.4cm}
}
\toprule
\textbf{Datensatz} &
\textbf{Bezugsquelle} &
\textbf{Zugangsverfahren} &
\textbf{Lizenz} \\
\midrule

SisFall
& SISTEMIC, Univ.\ de Antioquia (Kolumbien)
& Direktdownload; URL zeitweise nicht erreichbar
& CC~BY~4.0 \\

MobiFall / MobiAct
& BMI Lab, Hellenic Mediterranean University (Griechenland)
& Auf Anfrage per E-Mail; unterschriebener Nutzungsvertrag erforderlich
& Nur nicht-kommerzielle Forschung \\

UMAFall
& Figshare\slash Univ.\ de Málaga (DOI 10.6084/m9.figshare.4214283)
& Direktdownload ohne Registrierung
& Angabe der Autoren-Publikationen \\

KFall
& KAIST, Google Sites (\emph{sites.google.com/view/kfalldataset})
& Direktdownload ohne Registrierung
& Zitierpflicht \\

FallAllD
& IEEE DataPort (Open Access)
& Kostenfreier IEEE-Account (keine IEEE-Mitgliedschaft erforderlich)
& Open Access, Zitierpflicht \\

UniMiB SHAR
& SAL-Lab, Univ.\ Milano-Bicocca (Italien)
& Direktdownload; Spiegelkopie auf Kaggle verfügbar
& Nur nicht-kommerzielle Forschung \\

WEDA-FALL
& GitHub-Repository von J.~Marques (IST Lissabon)
& Direktdownload / \texttt{git clone}
& Zitierpflicht \\

\bottomrule
\end{tabularx}
\end{table}

Alle in dieser Übersicht genannten Datensätze sind für den akademischen
Gebrauch grundsätzlich zugänglich; die formalen Anforderungen
unterscheiden sich jedoch erheblich hinsichtlich Aufwand und
Bearbeitungszeit. Für die weitere Bearbeitung in Aufgabe~4 wurde
sichergestellt, dass die im nächsten Abschnitt als Primär- und
Sekundärdatensätze ausgewählten Ressourcen ohne administrative
Verzögerungen zur Verfügung stehen.

\section{Begründete Entscheidung}
\label{sec:dataset_decision}

Auf Basis der oben aufgeführten Kriterien wird \textbf{SisFall} als
primärer Datensatz für die Modellentwicklung in Aufgabe~4 ausgewählt.
Die Gründe für diese Wahl lassen sich wie folgt zusammenfassen:

\begin{enumerate}
\item Die \textbf{6-Achsen-Konfiguration} stimmt mit der in
      Kapitel~\ref{ch:comparison} begründeten Empfehlung überein.
\item Der \textbf{Messbereich} von $\pm 16$\,g bietet den in
      Abschnitt~\ref{sec:range} diskutierten Sicherheitspuffer gegen
      Clipping, insbesondere für periphere Körperpositionen.
\item Die \textbf{hohe Abtastrate} von 200\,Hz erlaubt ein Downsampling auf
      den in Abschnitt~\ref{sec:samplingrate} empfohlenen Bereich
      ($20$--$50$\,Hz) ohne Informationsverlust.
\item Die \textbf{große Probandenzahl} (38, davon 15 Ältere) und die
      Vielfalt an Sturz- und ADL-Typen (15 bzw.\ 19) ermöglichen eine
      belastbare probandenübergreifende Validierung.
\item Der \textbf{etablierte Benchmarkstatus} in der Literatur stellt sicher,
      dass die eigenen Ergebnisse unmittelbar mit publizierten Arbeiten
      verglichen werden können.
\end{enumerate}

Als ergänzender Datensatz bietet sich \textbf{KFall} an, da er eine
vergleichbare Sensorik, eine ausreichende Abtastrate und insbesondere die
zeitliche Pre-Impact-Annotation liefert, die für zukünftige Erweiterungen
des Modells relevant sein kann.

%------------------------------------------------------------------
\chapter{Verlauf der Modellentwicklung (Aufgabe~4 Zwischenstand)}
\label{ch:verlauf}

Die vorausgehenden Kapitel haben die theoretischen Grundlagen gelegt:
Kapitel~\ref{ch:comparison} begründet die Wahl einer 6-Achsen-IMU,
Kapitel~\ref{ch:properties} leitet den empfohlenen Arbeitspunkt ab
($\pm 16$\,g, $20$--$50$\,Hz, $\geq 12$~Bit) und
Kapitel~\ref{ch:datasets} wählt SisFall als primären Datensatz.
Das vorliegende Kapitel dokumentiert den bisherigen Verlauf der
praktischen Modellentwicklung im Rahmen von Aufgabe~4 und stellt die
ersten quantitativen Ergebnisse vor. Sämtliche Experimente wurden mit
Python~3.12 unter Verwendung von NumPy, SciPy, scikit-learn und
TensorFlow/Keras durchgeführt. Der vollständige Quellcode ist im
begleitenden Repository verfügbar.

\section{Datenaufbereitung}
\label{sec:preprocessing}

Als Trainingsdatensatz dient SisFall \cite{sucerquia2017sisfall} mit
4505 Aufzeichnungen (2707~ADL, 1798~simulierte Stürze) von
38~Probanden. Die Rohdaten liegen als kommaseparierte Textdateien mit
neun Spalten vor (zwei Beschleunigungssensoren und ein Gyroskop, jeweils
als ADC-Rohwerte). Für die Modellentwicklung werden ausschließlich die
Kanäle des ADXL345 ($\pm 16$\,g) und des ITG3200
($\pm 2000\,^{\circ}/\mathrm{s}$) verwendet, da sie der in
Kapitel~\ref{ch:properties} empfohlenen Spezifikation entsprechen und
zugleich mit dem für den Prototyp vorgesehenen MPU-6050 kompatibel sind.
Die ADC-Rohwerte werden anhand der Herstellerkonstanten in
physikalische Einheiten umgerechnet:

\begin{equation}
a\,[\mathrm{g}] = \mathrm{raw} \cdot \frac{2 \cdot 16}{2^{13}},
\qquad
\omega\,[^{\circ}/\mathrm{s}] = \mathrm{raw} \cdot \frac{2 \cdot 2000}{2^{16}}.
\label{eq:conversion}
\end{equation}

Anschließend wird das Signal von der ursprünglichen Abtastrate
($200$\,Hz) auf die Zielrate ($50$\,Hz) heruntergetastet. Die
Segmentierung erfolgt in gleitenden Fenstern der Länge $2$\,s mit
$50\,\%$ Überlappung. Bei Sturzaufzeichnungen wird das Fenster auf den
Zeitpunkt des maximalen Beschleunigungsbetrags (Impact) zentriert, um
sicherzustellen, dass das Klassifikationsziel tatsächlich die
sturzrelevante Signatur enthält und nicht einen beliebigen ruhigen
Abschnitt derselben Aufnahme. Die resultierenden Fenster haben die
Dimension $100 \times 6$ (100~Abtastwerte $\times$ 6~Kanäle).

\section{Modell~1: Schwellwert-Baseline}
\label{sec:threshold}

Als erste Referenz dient ein klassischer Schwellwertalgorithmus nach
Bourke \cite{bagala2012}: Ein Fenster wird genau dann als Sturz
klassifiziert, wenn der Betrag der Beschleunigung sowohl einen unteren
Freifall-Schwellwert von $0{,}6$\,g unterschreitet als auch einen
oberen Impact-Schwellwert von $2{,}0$\,g überschreitet. Dieser
Algorithmus benötigt weder Training noch Merkmalsextraktion und bildet
damit die einfachste realisierbare Baseline.

\textbf{Ergebnis:} Sensitivität $96{,}7\,\%$, Spezifität $76{,}8\,\%$,
Genauigkeit $77{,}5\,\%$. Die hohe Sensitivität bestätigt, dass das
Schwellwertverfahren die meisten Stürze erkennt. Die niedrige
Spezifität zeigt jedoch das in Kapitel~\ref{ch:comparison} diskutierte
Kernproblem: Sturzähnliche Alltagsaktivitäten wie schnelles Hinsetzen
lösen ebenfalls Alarm aus.

\section{Modell~2: Random Forest mit Magnitude-Features}
\label{sec:randomforest}

Um die Spezifität zu verbessern, werden aus jedem Fenster 18
rotationsinvariante Merkmale extrahiert, die ausschließlich auf dem
Betrag (\ac{SVM}) der Beschleunigung und der Winkelgeschwindigkeit
basieren. Zu den Merkmalen zählen unter anderem Mittelwert,
Standardabweichung, Maximum, Minimum, Spannweite, RMS-Wert, Energie,
Schiefe, Kurtosis, Jerk-Standardabweichung und Signal Magnitude Area.
Die bewusste Beschränkung auf Betragsgrößen macht die Merkmale
weitgehend unabhängig von der Orientierung des Sensors am Körper und
trägt damit zur Robustheit gegenüber dem in Abschnitt~\ref{sec:domain}
diskutierten Positionswechsel bei.

Ein Random-Forest-Klassifikator mit 200~Bäumen und balancierter
Klassengewichtung wird mittels subjektunabhängiger
Kreuzvalidierung (GroupKFold, 5~Folds) bewertet, wobei die
Gruppierungsvariable der Proband ist. Damit wird sichergestellt, dass
keine Daten desselben Probanden gleichzeitig in Training und Test
auftreten.

\textbf{Ergebnis:} Sensitivität $94{,}2\,\%$, Spezifität $99{,}6\,\%$,
Genauigkeit $99{,}4\,\%$. Die Spezifität steigt gegenüber der
Schwellwert-Baseline von $76{,}8\,\%$ auf $99{,}6\,\%$, was die Anzahl
falsch-positiver Alarme um den Faktor~$50$ reduziert. Die
Feature-Importance-Analyse des Random Forest bestätigt die
Schlussfolgerung aus Kapitel~\ref{ch:comparison}: Vier der fünf
wichtigsten Merkmale basieren auf dem Gyroskop
(\texttt{gyr\_std}, \texttt{gyr\_max}, \texttt{gyr\_energy},
\texttt{gyr\_rms}), was den Beitrag der sechsten Achse zur
Falsch-Positiv-Reduktion quantitativ belegt.

\section{Modell~3: 1D-CNN}
\label{sec:cnn}

Als drittes und leistungsfähigstes Modell wird ein eindimensionales
Convolutional Neural Network (1D-CNN) trainiert, das direkt auf den
rohen 6-Kanal-Fenstern ($100 \times 6$) arbeitet und keine
handgefertigte Merkmalsextraktion benötigt. Die Architektur umfasst
drei Faltungsblöcke mit aufsteigender Filterzahl (32, 64, 128),
jeweils gefolgt von Batch-Normalisierung und Max-Pooling, sowie einen
Global-Average-Pooling-Kopf und zwei vollvernetzte Schichten mit
Dropout ($p = 0{,}3$). Das Modell wird mit dem Adam-Optimierer
(Lernrate $10^{-3}$), binärer Kreuzentropie und klassenbalancierten
Gewichten über 30~Epochen trainiert.

Die Bewertung erfolgt auf einem subjektunabhängigen Testsplit
($25\,\%$ der Probanden, 10~Subjekte, insgesamt 13\,031~Fenster).

\textbf{Ergebnis:} Sensitivität $98{,}0\,\%$, Spezifität $99{,}7\,\%$,
Genauigkeit $99{,}6\,\%$. Im Vergleich zum Random Forest sinkt die
Anzahl übersehener Stürze von 105 auf 18 (False Negatives), während
die Anzahl der Fehlalarme von 215 auf 39 (False Positives) zurückgeht.
Das Modell umfasst 45\,601~Parameter (178\,KB) und konvergiert
innerhalb von 30~Epochen ohne Anzeichen von Overfitting, wie die
Trainingskurven in Abbildung~\ref{fig:training} belegen.

% Platzhalter für die Trainingskurven
% \begin{figure}[H]
%   \centering
%   \includegraphics[width=0.9\textwidth]{cnn_training.png}
%   \caption{Verlauf von Loss und Sensitivität über 30 Epochen (Training und Validation).}
%   \label{fig:training}
% \end{figure}

\section{Vergleich der drei Modelle}
\label{sec:comparison_models}

Tabelle~\ref{tab:results} fasst die Ergebnisse zusammen.

\begin{table}[H]
\centering
\caption{Vergleich der drei Sturzerkennungsmodelle auf SisFall (subjektunabhängig).}
\label{tab:results}
\small
\begin{tabularx}{\textwidth}{l C{2cm} C{2cm} C{2cm} C{1.5cm} C{1.5cm}}
\toprule
\textbf{Modell} & \textbf{Sensitivität} & \textbf{Spezifität} & \textbf{Genauigkeit} & \textbf{FP} & \textbf{FN} \\
\midrule
Schwellwert (Bourke)   & 96,7\,\% & 76,8\,\% & 77,5\,\% & 11\,503 & 59 \\
Random Forest          & 94,2\,\% & 99,6\,\% & 99,4\,\% & 215     & 105 \\
\textbf{1D-CNN}        & \textbf{98,0\,\%} & \textbf{99,7\,\%} & \textbf{99,6\,\%} & \textbf{39} & \textbf{18} \\
\bottomrule
\end{tabularx}
\end{table}

Der stufenweise Aufbau zeigt, dass jedes nachfolgende Modell den
vorhergehenden Schwachpunkt adressiert: Der Random Forest behebt die
niedrige Spezifität der Schwellwertmethode, und das 1D-CNN verbessert
zusätzlich die Sensitivität bei gleichzeitig nochmals reduzierter
Fehlalarmrate.

\section{TFLite-Export für den Mikrocontroller}
\label{sec:tflite}

Für den Einsatz auf dem ESP32-C6 wird das trainierte Keras-Modell
mittels des TensorFlow-Lite-Konverters in ein vollständig
int8-quantisiertes Modell überführt. Die Quantisierung verwendet einen
repräsentativen Kalibrierungsdatensatz aus 200~zufällig gewählten
Trainingsfenstern. Das resultierende \texttt{.tflite}-Modell hat eine
Größe von lediglich \textbf{64\,KB} und passt damit problemlos in den
4\,MB großen Flash-Speicher des ESP32-C6. Die Inferenzzeit für ein
einzelnes Fenster wird auf dem Zielprozessor im Bereich weniger
Millisekunden erwartet, sodass eine Echtzeitklassifikation bei der
gewählten Abtastrate ($50$\,Hz, Fenster alle $1$\,s) gewährleistet ist.

\section{Domänenlücke und Mitigationsstrategien}
\label{sec:domain}

Der SisFall-Datensatz wurde mit einer am Gürtel (Hüfte) getragenen
Sensoreinheit aufgenommen, während das Zielsystem eine
schuhsohlenintegrierte Anbringung vorsieht. Diese Positionsänderung
stellt die wesentliche Domänenlücke dar, da sich die Kinematik der
Fußbewegung grundlegend von der des Rumpfes unterscheidet: Beim Gehen
erzeugt der Fersenaufsatz (Heel Strike) Beschleunigungsspitzen, die
denen eines leichten Stoßereignisses ähneln und somit die
Falsch-Positiv-Rate erhöhen können.

Um die Übertragbarkeit zu verbessern, wurden bereits in der aktuellen
Pipeline mehrere Maßnahmen umgesetzt:

\begin{enumerate}
\item \textbf{Rotationsinvariante Merkmale:} Die Random-Forest-Features
      basieren ausschließlich auf Betragsgrößen (\ac{SVM}), die
      unabhängig von der Orientierung des Sensors sind.
\item \textbf{Per-Window-Normalisierung:} Jedes Fenster wird vor der
      Eingabe in das CNN auf Mittelwert~0 und Standardabweichung~1
      normiert, wodurch positionsabhängige Gleichanteile (z.\,B.\
      unterschiedliche Gravitationsprojektion) eliminiert werden.
\item \textbf{Subjektunabhängige Validierung:} Die Aufteilung nach
      Probanden stellt sicher, dass die berichteten Kennzahlen nicht
      durch probandenspezifische Muster überschätzt werden.
\end{enumerate}

Für die verbleibende Domänenlücke ist ein \emph{Pilot-Fußdatensatz}
geplant: Mit dem für den Prototyp vorgesehenen MPU-6050 auf dem
ESP32-C6 werden Daten direkt am Fuß aufgezeichnet (Gehen, Laufen,
Treppensteigen, Hinsetzen sowie simulierte Stürze auf eine Matte).
Anschließend werden die letzten Schichten des trainierten CNN mittels
Transfer Learning auf diesen Datensatz feinabgestimmt (Fine-Tuning).
Erste Arbeiten mit dem UMAFall-Knöcheldatensatz
\cite{umafall2017}, der die nächstgelegene öffentlich verfügbare
Sensorposition bietet, sind als Zwischenschritt vorgesehen.

\section{Nächste Schritte}
\label{sec:nextsteps}

Die folgenden Arbeitsschritte sind für den verbleibenden
Projektzeitraum geplant:

\begin{enumerate}
\item \textbf{Ablationsstudie zur Abtastrate:} Wiederholung des
      CNN-Trainings bei $10$, $20$ und $50$\,Hz, um die in
      Kapitel~\ref{ch:properties} diskutierte Empfehlung experimentell
      zu überprüfen.
\item \textbf{Kreuzvalidierung auf KFall:} Training auf SisFall,
      Test auf KFall \cite{kfall2021} und umgekehrt, um die
      Generalisierbarkeit über Datensatzgrenzen hinweg zu bewerten.
\item \textbf{Hardware-Integration:} Aufbau des ESP32-C6 + MPU-6050
      Prototyps, Portierung des int8-TFLite-Modells mittels
      \texttt{esp-tflite-micro} oder Edge Impulse und Messung der
      On-Device-Inferenzzeit sowie des Energieverbrauchs.
\item \textbf{Pilotdatenerhebung am Fuß:} Aufzeichnung eines kleinen
      Fußdatensatzes zur Validierung und zum Fine-Tuning.
\item \textbf{Long-Lie-Erkennung:} Implementierung einer zweiten
      Erkennungsstufe, die nach einem erkannten Sturz die
      Bewegungslosigkeit überwacht und bei anhaltender Liegeposition
      ($> 60$\,s) einen Notfallalarm auslöst.
\end{enumerate}

%------------------------------------------------------------------
%   BIBLIOGRAPHY
%------------------------------------------------------------------
\newpage
\bibliography{references}

%------------------------------------------------------------------
%   VERSICHERUNG (Selbstständigkeitserklärung)
%------------------------------------------------------------------
\newpage
\thispagestyle{empty}
\vspace*{2cm}
\section*{Versicherung}
Ich versichere, dass ich diese Arbeit selbstständig verfasst und keine anderen
als die angegebenen Hilfsmittel benutzt habe. Die den benutzten Hilfsmitteln
wörtlich oder inhaltlich entnommenen Stellen habe ich unter Quellenangaben
kenntlich gemacht. Die Arbeit hat in gleicher oder ähnlicher Form noch keiner
anderen Prüfungsbehörde vorgelegen.

\vspace{1cm}
\noindent I declare that I have authored this thesis independently, that I have
not used other than the declared sources / resources and that I have explicitly
marked all material which has been quoted either literally or by content from
the used sources.

\vspace{3cm}
\noindent
\rule{6cm}{0.4pt}\hfill\rule{6cm}{0.4pt}\\
\hspace*{1cm}Place, Date\hfill Aly Elbermawy\hspace*{2cm}

\end{document}